\documentclass{article}

\PassOptionsToPackage{numbers, sort, compress}{natbib}
\usepackage[main,final]{neurips_2025}

\usepackage[T1]{fontenc}
\usepackage[utf8]{inputenc}
\usepackage{hyperref}
\usepackage{url}
\usepackage{booktabs}
\usepackage{nicefrac}
\usepackage{microtype}
\usepackage{xcolor}

%%%

\usepackage{subcaption}
\usepackage{graphicx}
\usepackage{multirow}
\usepackage{amsmath,amssymb,amsfonts}
\usepackage{amsthm}
\usepackage{mathrsfs}
\usepackage{xcolor}
\usepackage{textcomp}
\usepackage{manyfoot}
\usepackage{algorithm}
\usepackage{algorithmicx}
\usepackage{algpseudocode}
\usepackage{listings}

\newtheorem{theorem}{Theorem} % continuous numbers
%%\newtheorem{theorem}{Theorem}[section] % sectionwise numbers
%% optional argument [theorem] produces theorem numbering sequence instead of independent numbers for Proposition
\newtheorem{proposition}[theorem]{Proposition}% 
\newtheorem{lemma}{Lemma}% 
%%\newtheorem{proposition}{Proposition} % to get separate numbers for theorem and proposition etc.

\newtheorem{example}{Example}
\newtheorem{remark}{Remark}

\newtheorem{definition}{Definition}
\newtheorem{assumption}{Assumption}

%%%


\title{Calibrated Perplexity for Robust AI-Generated Text Detection Across Languages and Genres}


% The \author macro works with any number of authors. There are two commands
% used to separate the names and addresses of multiple authors: \And and \AND.
%
% Using \And between authors leaves it to LaTeX to determine where to break the
% lines. Using \AND forces a line break at that point. So, if LaTeX puts 3 of 4
% authors names on the first line, and the last on the second line, try using
% \AND instead of \And before the third author name.


\author{
  Name Surname\\
  MIPT\\
  Moscow, Russia\\
  \texttt{email@phystech.edu}\\
  \And
  Name Surname\\
  MIPT\\
  Moscow, Russia\\
  \texttt{email@phystech.edu}\\
}


\begin{document}


\maketitle

\begin{abstract}
    This paper invistigates the problem of robust AI-generated text detection under distribution shift. While simple perplexity is a simple solution for such detection, its oppurtunities lose significantly on out-of-domain data. We propose a novel method - calibrated perplexity - which normalizes raw perplexity scores using domain-specific language model baselines. This approach constructs a stable feature space less sensitive to text length, genre, and language variations. Experiments across multiple domains and languages are going to demonstrate improved in-domain and out-of-domain detection accuracy compared to standard perplexity-based methods.
\end{abstract}

\textbf{Keywords:} Calibrated perplexity
Cross-domain detection
Robustness
Machine-generated text
Text classification
Multilingual detection
Genre adaptation
Statistical methods
Language modeling

\section{Introduction}\label{sec:intro}

Fast improvement and widespread accessibility of Large Language Models (LLMs) have destroyed the lines between human-written and machine-generated text. While this technology offers benefits, it also presents significant disadvantages, including academic dishonesty, disinformation, and the fall of trust in digital content \cite{xu2024detecting}, \cite{hans2024spotting}. This has created an inevitable need for robust methods to detect AI-generated content.

Current detectors often rely on fine-tuned classifiers or statistical signatures taken from the text \cite{gehrmann2019gltr}. Useful and computationally efficient feature is perplexity, a measure of how "surprising" a text is to a language model. The assumption is that LLMs, which are trained to maximize likelihood, will generate text with lower perplexity compared to human writing. As it is shown, a significant change comes from distribution shift. The perplexity of a text is not an absolute value; it is highly dependent on the domain (e.g., news, scientific articles, social media) and language it was trained on. A model fine-tuned for one genre often fails on another, causing unreliable detection in real world, where the target domain is unknown \cite{bakhteev2022automatic}.

Recent works have attempted to solve this by developing zero-shot detectors that are more robust to unseen prompts or models. For instance, Binoculars \cite{hans2024spotting} uses a ratio of scores from two different LLMs for calibration. Similarly, research on multilingual and cross-domain detection highlights the necessity of adapting methods to diverse data distributions \cite{gritsay2023artificially}. Despite these advances, a method for calibrating perplexity itself to make it robust across different genres and languages remains still not-explored.

In this paper, we face the challenge of domain shift in AI-generated text detection. We propose a novel method called calibrated perplexity, which normalizes perplexity scores against domain-specific baselines. Our approach includeses training KenLM n-gram models on text from multiple domains and using these models to compute a relative, rather than absolute, measure of text predictability. The core novelty of our work lies in transforming a domain-dependent score into a stable feature for classification. To check our method, we conduct experiments across multiple domains and languages, comparing in-domain and out-of-domain detection.


\section{Problem statement}

Let $\mathbf{W}$ be our alphabet, tokens consists from characters from that alphabet.

Let $$\mathbb{D} = \Bigl\{\Big[t_j\Big]_{j=1}^n \Hquad|\Hquad t_j \in \mathbf{W}, n \in \mathbb{N}\Bigr\}$$ be the space of our documents.

We have a set of $N$ documents
$$\mathbf{D} = \bigcup_{i=1}^{N}D^i, D^i \in \mathbb{D}$$

%Each document  $D^i$ consists of $d_i$ tokens $t_1^i, t_2^i, ..., t_{d_i}^i$, where $t_j^i \in \mathbf{W}$ .
We have $K$ classes of our text-generative models, that generated fragments of text in $\mathbf{D}$.

We have $K + 1$ labels, that we use for multiclass classification, where $0$ is the label of human-written fragment, $\{1...K\}$ are the  labels that represent corresponding language model, let $\mathbf{C}$ be our set of labels.

Our model $\mathbf{\phi}$ is a superposition of two mappings, $\mathbf{f}$ and $\mathbf{g}$. Mapping $\mathbf{f}$ is responsible for text segmentation. Mapping $\mathbf{g}$ is responsible for classifying obtained fragments.

$$\mathbf{\phi} : \mathbf{g} \circ \mathbf{f},$$


$$\mathbf{\phi}: \mathbb{D} \rightarrow \mathbb{T}, \qquad \mathbb{T} = \Bigl\{\Big[t_{s_j}, t_{f_j}, C_j\Big]_{j = 1}^{J} \Hquad|\Hquad t_{s_j} = t_{f_{j - 1}},\Hquad s_j \in \mathbb{N}_0,\Hquad f_j \in \mathbb{N}, \Hquad C_j \in \mathbf{C}\Bigr\},$$

where $J$ is a number of fragments in segmentation of the document, $t_{s_j}$ is starting index of the $j$-th fragment,  $t_{f_j}$ is finishing index of the $j$-th fragment,  $C_{j}$ is class of the $j$-th fragment.


To estimate the statistical properties of text fragments, we use a character-aware n-gram language model, KenLM \cite{chen1996empirical}. Let a fragment of text be represented as a sequence of tokens $[t_a, t_{a+1}, ..., t_b]$, where $t_i \in \mathbf{W}$. KenLM provides a function $\mathcal{K}$ that computes the log-probability of this sequence under an n-gram model trained on a specific corpus.

Formally, we define a language model $\mathcal{M}_{\mathcal{T}}$ trained on a text corpus $\mathcal{T} \subset \mathbb{D}$. For a given text fragment $F = [t_{a}, ..., t_{b}]$, the model estimates its probability as the product of conditional probabilities of each token given its $n-1$ predecessors:
\begin{equation}
    P_{\mathcal{M}_{\mathcal{T}}}(F) = \prod_{i=a}^{b} P_{\mathcal{M}_{\mathcal{T}}}(t_i \mid t_{i-1}, ..., t_{i-n+1}).
\end{equation}
KenLM efficiently implements this estimation using modified Kneser-Ney smoothing, which is particularly effective for handling sparse n-gram occurrences.

The **perplexity** of a fragment $F$ is then defined as the inverse geometric mean of the token probabilities:
\begin{equation}
    \text{PPL}_{\mathcal{M}_{\mathcal{T}}}(F) = P_{\mathcal{M}_{\mathcal{T}}}(F)^{-\frac{1}{|F|}} = 2^{-\frac{1}{|F|}\log_2 P_{\mathcal{M}_{\mathcal{T}}}(F)},
\end{equation}
where $|F| = b-a+1$ is the length of the fragment.

This metric quantifies how "surprising" the fragment is to the model $\mathcal{M}_{\mathcal{T}}$; lower perplexity indicates that the fragment is more predictable based on the training corpus $\mathcal{T}$. In our work, we train multiple KenLM models $\{\mathcal{M}_{\mathcal{T}_d}\}$ on different domains $\mathcal{T}_d$ (e.g., news, science, social media). For a given fragment $F$, we compute a set of raw perplexities $\{\text{PPL}_{\mathcal{M}_{\mathcal{T}_d}}(F)\}$. The subsequent mapping $\mathbf{g}$ then uses these values, along with the segmentation information from $\mathbf{f}$, to perform the final classification.

\textbf{Contributions.} Our contributions can be summarized as follows:
\begin{itemize}
    \item We present...
    \item We demonstrate the validity of our theoretical results through empirical studies...
    \item We highlight the implications of our findings for...
\end{itemize}

\textbf{Outline.} The rest of the paper is organized as follows...

\section{Related Work}\label{sec:rw}

\textbf{Topic \#1.}
TODO

\textbf{Topic \#2.}
TODO

\section{Preliminaries}\label{sec:prelim}

\subsection{General notation}

In this section, we introduce the general notation used in the rest of the paper and the basic assumptions. 

\subsection{Assumptions} 

TODO

\section{Method}\label{sec:method}

\section{Experiments}\label{sec:exp}

To verify the theoretical estimates obtained, we conducted a detailed empirical study...

\section{Discussion}\label{sec:disc}

TODO

\section{Conclusion}\label{sec:concl}

TODO


%%%%%%%%%%%%%%%%%%%%%%%%%%%%%%%%%%%%%%%%%%%%%%%%%%%%%%%%%%%%

\bibliographystyle{unsrtnat}
\bibliography{references}

%%%%%%%%%%%%%%%%%%%%%%%%%%%%%%%%%%%%%%%%%%%%%%%%%%%%%%%%%%%%

\newpage
\appendix
\section{Appendix / supplemental material}\label{app}

\subsection{Additional experiments / Proofs of Theorems}\label{app:exp}

TODO

\end{document}
